\documentclass[tikz, border=3mm]{standalone}
\usepackage[spanish]{babel}
\usepackage[utf8]{inputenc}
\usepackage[T1]{fontenc}
\usepackage{xcolor}

\begin{document}
\shorthandoff{>}
\begin{tikzpicture}[every node/.style={draw, rounded corners, align=center, font=\footnotesize}, >=stealth]
 % Define colors
 \definecolor{netpurple}{HTML}{7C3AED}
 \definecolor{serverblue}{HTML}{4C5CC5}
 \definecolor{laptopgreen}{HTML}{00B299}
 \definecolor{telegramblue}{HTML}{0088CC}

 % --- LEVEL 0: WAN / ROUTER ---
 \node[fill=netpurple!20, text width=2.6cm, minimum height=1.2cm] (router) at (0,2.5) {\textbf{TL-MR6400}\\{\footnotesize 192.168.1.1}\\{\tiny Sin SIM / sin WAN}};

 % Internet connection status
 \node[draw=none, font=\tiny, red!70] at (0,4.1) {\textbf{Sin Internet}};
 \draw[dashed, red!40] (-1.2,3.8) -- (1.2,3.8);
 \node[draw=none, font=\tiny, red!40] at (0,3.5) {(ranura SIM rota)};

 % --- LEVEL 1: IP NETWORK DEVICES ---
 \node[fill=laptopgreen!20, text width=2.6cm, minimum height=0.9cm] (laptop) at (-4.5,0) {\textbf{Portátil}\\{\footnotesize 192.168.1.101}\\{\tiny Desarrollo + SSH}};

 \node[fill=serverblue!20, text width=2.8cm, minimum height=0.9cm] (server) at (0,0) {\textbf{Torre Ryzen}\\{\footnotesize 192.168.1.50}\\{\tiny Ubuntu + Ollama + OpenClaw}};

 \node[fill=gray!20, text width=2.6cm, minimum height=0.9cm] (tuya) at (4.5,0) {\textbf{Enchufe Tuya}\\{\footnotesize 192.168.1.100}};

 % Physical Network Backbone Bus
 \coordinate (bus) at (0,1.2);
 \draw[thick] (router.south) -- (bus);
 \draw[thick, densely dashed] (bus) -| (laptop.north);
 \draw[thick] (bus) -- (server.north);
 \draw[thick] (bus) -| (tuya.north);

 % Dedicated Development Link
 \draw[thick, laptopgreen, <->] (laptop.east) -- (server.west) node[midway, above, font=\tiny, draw=none, text=black] {};

 % Server to Tuya Plug IP Control Path (Routed via Network Layer)
 \draw[thick, dashed, ->, serverblue] ([xshift=0.2cm]server.north) -- ++(0,0.5) -| ([xshift=-0.2cm]tuya.north) node[pos=0.35, above, font=\tiny, text=black] {Control IP};

 % --- LEVEL 2: CONTROL LAYER & UNCONNECTED SYSTEMS ---
 \node[fill=telegramblue!20, text width=2.6cm, minimum height=0.7cm] (telegram) at (-4.5,-2.5) {\textbf{Canal Telegram}\\{\tiny Comunicación Bot}};

 \node[fill=gray!20, text width=2.6cm, minimum height=0.7cm] (esp32) at (0,-2.5) {\textbf{ESP32}\\{\footnotesize Nodo Control}};

 % Isolated HVAC system (Pending)
 \node[fill=gray!15, text width=2.6cm, minimum height=0.7cm, dashed] (hvac) at (4.5,-2.5) {\textbf{HVAC}\\{\footnotesize (pendiente)}};

 % Telegram Bot to Server communication path
 \draw[thick, telegramblue, <->] (telegram.east) -- ++(1,0) |- ([xshift=-0.3cm]server.south) node[pos=0.25, above, font=\tiny, draw=none, text=black] {Bot API};

 % Server to ESP32 Management Link
 \draw[thick, serverblue, ->] (server.south) -- (esp32.north) node[midway, right, font=\tiny, draw=none, text=black] {Gestión};

 % --- LEVEL 3: BLE PERIPHERALS ---
 \node[fill=gray!10, text width=2.6cm, minimum height=0.7cm] (xiaomi) at (-2,-4.8) {\textbf{Sensor Xiaomi}\\{\footnotesize Métrica BLE}};

 \node[fill=gray!20, text width=2.6cm, minimum height=0.7cm] (fan) at (2,-4.8) {\textbf{FanLamp F8808}\\{\footnotesize BLE ads}};

 % Pure BLE Control Lines originating from ESP32
 \draw[thick, densely dotted, ->] (esp32.south) -- ++(0,-0.6) -| (xiaomi.north) node[pos=0.75, above, font=\tiny, draw=none, text=black] {BLE};
 \draw[thick, densely dotted, ->] (esp32.south) -- ++(0,-0.6) -| (fan.north) node[pos=0.75, above, font=\tiny, draw=none, text=black] {BLE};

 % Subnet Boundary Indicator
 \draw[dashed, thick, netpurple!30] (-6,-5.8) -- (6,-5.8);
 \node[draw=none, font=\tiny, netpurple!60] at (4.5,-5.5) {segmento 192.168.1.0/24};

\end{tikzpicture}
\end{document}
